% TOP_PROBLEM: {"id": 2073, "title": "Erdős Problem #359", "category_id": 1, "category": {"id": 1, "name": "number_theory", "display_name": "Number Theory", "description": "Properties of integers, prime numbers, Diophantine equations.", "slug": "number-theory", "order_index": 1, "created_at": "2026-07-31T15:26:25.670Z"}, "status": "open"}
% FIRST_POSTED: null
% ENTRY_KIND: statement_only
% Imported from: unsolvedmath_additional_500.tex; UnsolvedMath ID 2073
% !TeX program = lualatex
% Compile twice: lualatex 2073.tex
% Self-contained: no external bibliography, images, or \input files.
% Numeric section labels and visible numbers are original UnsolvedMath IDs.
\documentclass[11pt,a4paper]{article}
\usepackage[margin=24mm,headheight=14pt]{geometry}
\usepackage{fontspec}
\setmainfont{Latin Modern Roman}
\setsansfont{Latin Modern Sans}
\setmonofont{Latin Modern Mono}
\usepackage{amsmath,amssymb,mathtools}
\usepackage{unicode-math}
\setmathfont{Latin Modern Math}
\usepackage{microtype}
\usepackage{enumitem,xurl,needspace}
\usepackage[dvipsnames]{xcolor}
\usepackage{hyperref}
\hypersetup{unicode=true,colorlinks=true,linkcolor=MidnightBlue,urlcolor=MidnightBlue,
 pdftitle={UnsolvedMath Problem 2073},pdfauthor={Source-annotated compilation},bookmarksnumbered=true}
\usepackage{bookmark,tocloft,titlesec,fancyhdr}
\setlength{\cftsecnumwidth}{4.2em}
\cftsetpnumwidth{2.5em}
\cftsetrmarg{4em}
\titleformat{\section}{\Large\bfseries\raggedright}{\thesection}{0.7em}{}
\pagestyle{fancy}\fancyhf{}
\fancyhead[L]{\small\sffamily Additional UnsolvedMath statements}
\fancyhead[R]{\small Original catalogue IDs}
\fancyfoot[C]{\thepage}
\setlength{\parindent}{0pt}
\setlength{\parskip}{5pt plus 1pt}
\setlength{\emergencystretch}{3em}
\setcounter{tocdepth}{1}\setcounter{secnumdepth}{1}
\setlist[itemize]{leftmargin=3em,itemsep=4pt,topsep=4pt,parsep=0pt}
\allowdisplaybreaks
\newcommand{\N}{\mathbb{N}}
\newcommand{\Z}{\mathbb{Z}}
\newcommand{\Q}{\mathbb{Q}}
\newcommand{\R}{\mathbb{R}}
\newcommand{\C}{\mathbb{C}}
\newcommand{\F}{\mathbb{F}}
\newcommand{\A}{\mathbb{A}}
\newcommand{\PP}{\mathbb{P}}
\newcommand{\E}{\mathbb{E}}
\newcommand{\eps}{\varepsilon}
\newcommand{\dd}{\,\mathrm d}
\newcommand{\sref}[2]{\hyperlink{source:#1:#2}{[S#2]}}
\newif\ifshowcatalogue
\showcataloguetrue
\newcommand{\cataloguescope}[1]{\ifshowcatalogue
 \paragraph{Statement recorded in the supplied source.}
 #1\fi}
\begin{document}
% UnsolvedMath ID 2073; source catalogue code EP-359

\setcounter{section}{2072}

\section{Density of a Greedy Consecutive-Sum-Avoiding Sequence}\label{2073}

{\small\sffamily UnsolvedMath ID 2073 · Number Theory}

\subsection{Definitions and mathematical statement}

\paragraph{Shared notation.}Unless a section says otherwise, $\N=\{1,2,\ldots\}$,
$\Z$ is the ring of integers, $\Q,\R,\C$ are the rational, real, and complex
fields, $[N]=\{1,\ldots,\lfloor N\rfloor\}$, and $\log$ is the natural
logarithm. For real functions of a parameter tending to infinity,
$f\ll g$ and $f=O(g)$ mean $|f|\leq Cg$ eventually for some fixed $C$;
$f\gg g$ reverses this comparison; $f=o(g)$ means $f/g\to0$;
$f\sim g$ means $f/g\to1$; and $f\asymp g$ means both order bounds.
A subscript on $O$ or $\ll$ specifies allowed constant dependence.
The expression $n^{o(1)}$ in an upper bound means growth smaller than
$n^\epsilon$ for each fixed $\epsilon>0$. Local definitions govern any
exception, finite-field convention, or use of the same symbol for another
object. An asymptotic claim is not automatically an effective algorithm.



The increasing ordering makes the next choice the least integer greater than the previous term that is not the sum of a nonempty consecutive block of earlier terms. The density question concerns the set of values, with the initial seed held fixed.

For an integer set $A$, its natural density is $\lim_{X\to\infty}|A\cap[1,X]|/X$, when this limit exists. Upper and lower densities use limit superior and limit inferior. A request for density is not automatically satisfied by positive upper density alone. \sref{2073}{1} \sref{2073}{2}

\paragraph{Statement recorded in the supplied source.}
Let $a_1<a_2<\cdots$ be an infinite sequence of integers such that $a_1=n$ and $a_{i+1}$ is the least integer which is not a sum of consecutive earlier $a_j$s. What can be said about the density of this sequence?
In particular, in the case $n=1$, can one prove that $a_k/k\to \infty$ and $a_k/k^{1+c}\to 0$ for any $c>0$? \sref{2073}{1}

\subsection{Short English statement}

What is the density of the sequence that greedily avoids sums of earlier consecutive terms, especially when it begins with one?

\subsection{Sources}

{\small

\begin{itemize}

\item[\hypertarget{source:2073:1}{S1}] ULAM.AI, \emph{UnsolvedMath}, supplied \texttt{problems.json}, record 2073 (EP-359), statement and background. The embedded literature-review date is 2026-08-17; that is the archive’s review date, not a fresh certification. Dataset landing page: \url{https://huggingface.co/datasets/ulamai/UnsolvedMath}.

\item[\hypertarget{source:2073:2}{S2}] T. F. Bloom, \emph{Erdős Problems}, Problem 359, \url{https://www.erdosproblems.com/359}. Reference imported from the supplied record; the live problem page was not independently checked for this compilation.

\item[\hypertarget{source:2073:3}{S3}] Andrews, George E., Research Problems: Mac Mahon's Prime Numbers of Measurement. Amer. Math. Monthly (1975), 922-923. Bibliographic information imported from the supplied record.

\item[\hypertarget{source:2073:4}{S4}] Porubsk\'y, \v S., On {M}ac{M}ahon's segmented numbers and related sequences. Nieuw Arch. Wisk. (3) (1977), 403--408. Bibliographic information imported from the supplied record.

\end{itemize}

}

\label{end:2073}
\end{document}
